\section{A bounded-exponent determinant construction and a strict test case}

\subsection{O-014.a: degree splitting with exponent bounds}
Let $b\geq1$, $r\geq2$, and choose positive integer caps
\[
             c=(c_2,\ldots,c_r),\qquad c_2=b,
                    \quad c_i\leq b\quad(i\geq3).
\]
Write $u=\lfloor(b-1)/2\rfloor$, $v=\lceil(b-1)/2\rceil$,
and define
\[
 N_j(c)=\#\{\alpha\in\mathbb Z_{\geq0}^{r-1}:
                         \alpha\leq c,\ |\alpha|\leq j\},
 \qquad D(b,c)=N_u(c)+N_v(c)-1.                         \tag{O-014.1}
\]
Equivalently,
\[
 N_j(c)=[z^j]\frac{1}{1-z}\prod_{i=2}^r
                                      (1+z+\cdots+z^{c_i}).
\]
For every degree-$b$ form
\[
 f(L,z_2,\ldots,z_r)=
       \sum_{|\gamma|\leq b,\,\gamma\leq c}
                          a_\gamma z^\gamma L^{b-|\gamma|},
\]
there is a homogeneous linear matrix of size $D=D(b,c)$
whose determinant equals $L^{D-b}f$. Its entries are linear
forms in the form variables, with integral polynomial dependence
on the coefficient parameters.
In particular, every $L^{n-b}f$ in this bounded coefficient
space lies in $\Omega_n$ if $n\geq D$ and $r\leq n^2$.

\begin{proof}
Use the two trees of O-012.a, retaining only multi-indices
$\alpha\leq c$ in the first tree of height $u$ and the second
tree of height $v$. The retained sets are closed under taking
parents, so every root path remains present. They have $N_u(c)$
and $N_v(c)$ vertices; merge only their roots.
For a coefficient with $|\gamma|<b$, split its sorted word
into $\alpha+\beta=\gamma$ with $|\alpha|\leq u$ and
$|\beta|\leq v$, and insert the coefficient edge labeled
$a_\gamma L$. Both endpoints satisfy the caps because they
are coordinatewise bounded by $\gamma$. For $|\gamma|=b$,
remove one variable $z_i$, split $\gamma-e_i$ in the same
way, and label the coefficient edge by $a_\gamma z_i$.
Again both endpoints and every path are retained.

All tree edges carry the negative of their variable, all
noncentral unused vertices carry $L$-loops, and coincident
coefficient-edge labels are added. Every directed non-loop
cycle passes through the central vertex and exactly one
coefficient edge. The determinant signs cancel the tree-edge
signs just as in O-012.a, giving one term
$a_\gamma z^\gamma L^{D-|\gamma|}$ for each coefficient.
Their sum is $L^{D-b}f$. Since $c_2=b$, the first variable
alone supplies paths of full heights $u$ and $v$, so
$D\geq u+v+1=b$. No negative power of $L$ occurs.
Adding an $L I_{n-D}$ block gives the larger determinant.
\end{proof}

For a partition $\lambda$, use $b=\lambda_2$ and
$c=(\lambda_2,\ldots,\lambda_r)$. F-015 proves that this
whole restricted coefficient space is sufficient to preserve
every highest-weight copy, so that $n\geq D(b,c)$ implies
$K_n(\lambda)=a_\lambda(d[n])$. The argument uses the
second-row stabilization F-014 and the invertible diagonal
padding pullback of O-012.b. The size in (O-014.1) is at
most the unbounded size $B(b,r)$, and can be strictly smaller.

\subsection{O-014.b: a genuine positive-permanent example beyond the earlier gates}

Here is an exact test case demonstrating that the bounded
coefficient refinement can add an exclusion, including when
the ambient multiplicity is larger than one. Set
\[
 m=256,\quad r=m^2=65536,\quad b=32,\quad
 c=(32,2^{65534}),\quad n=D(32,c),\quad d=n.
\]
The exact integer is
\[
n=N_{15}(c)+N_{16}(c)-1.
                                                               \tag{O-014.2}
\]
Take the partition
\[
                 \lambda=(n^2-131100,32,2^{65534}).     \tag{O-014.3}
\]
Then
\[
              4\leq q^m_\lambda(d[n])\leq K_n(\lambda)
                     =a_\lambda(d[n]).
                                                               \tag{O-014.4}
\]
The comparison of every copy follows from F-015, not from
the explicit one-copy construction below.

To reproduce (O-014.2), put $R=65534$. Since $u=15$, $v=16$
and the first cap is inactive in these degrees,
\[
 N_j(c)=\sum_{k=0}^j(j-k+1)
          \sum_{t=0}^{\lfloor k/2\rfloor}
                         \binom Rt\binom{R-t}{k-2t}.
                                                               \tag{O-014.5}
\]
The inner sum selects the tail variables used twice and once;
the factor $j-k+1$ counts the possible first-coordinate
exponents in a monomial of total degree at most $j$.
Thus (O-014.5), followed by $n=N_{15}+N_{16}-1$, is a
finite integer certificate independent of decimal rounding.
The complete decimal integer is retained in
\texttt{boxed\_degree\_split\_certificate.json}.

\paragraph{Positive permanent multiplicity.}
For the internally padded permanent, choose the nondegenerate
Hessian point from O-011 and an orthogonal basis beginning
with its nonisotropic radial direction. Every leading Hessian
minor in that basis is nonzero, including $W_2$ and $W_r$;
the form itself is nonzero in the first direction. Therefore
\[
                    W_r(F)W_2(F)^{15}F(e_1)^{d-r-30}
                                                               \tag{O-014.6}
\]
is a nonzero degree-$d$ highest-weight function on the
permanent orbit. Its second row is $2+30=32$, its remaining
nonfirst rows are $2$, and its first row is determined by
total size $nd$. Its weight is exactly (O-014.3), proving
$q^m_\lambda>0$. The stronger bound $q^m_\lambda\geq4$
follows from M-012.3, whose hypotheses hold because $b=32$
and $d\geq r+30$. It retains four independent binary seed
classes after multiplication by one common nonzero Hessian
factor on the integral permanent coordinate ring.

For the ambient lower bound, use F-013's four-dimensional
highest-weight space at
\[
                    (r(n-2)-4,8,2^{r-2})
\]
in outer degree $r$. Multiplication by the nonzero polynomial
$W_2^{12}F(e_1)^{d-r-24}$ sends all four vectors to the
highest-weight space at (O-014.3). The ambient polynomial
ring is a domain, so multiplication by this fixed polynomial
is injective. Thus $a_\lambda\geq4$. No value for the actual
exact permanent multiplicity is inferred from this ambient calculation.

\paragraph{Comparison with previously proved sufficient conditions.}
Exact integer arithmetic in the accompanying certificate gives
\[
                   m^{25}\leq n<2^m-1,
\]
so the BIP range holds and the full Grenet containment gate
F-009 does not apply. The unbounded size is
\[
B(32,65536)=\binom{65550}{15}+\binom{65551}{16}-1>n,
\]
so the earlier second-row determinant criterion also fails.
For any split of the $r-1=65535$ tail variables, some part
has at least $32768$ variables. Since all caps are at least
two, its monomial count through degree $m$ is at least
$\binom{32768}{256}$, by choosing squarefree monomials of
degree $m$. Consequently
\[
        T_m(\lambda)\geq\binom{32768}{256}>n.
\]
The analogous degree-$b$ variable-splitting bound is at least
$\binom{32768}{32}>n$. The uniform row cutoff does not apply,
since $r>51$, and the second row is greater than $25$.
Since $d=n$, the quadratic Chow bound $n\geq d(d-1)$
and the adaptive bound $n\geq B_d(32)=32d-497$ also fail.
The coarse padding requirements hold: $r=m^2$ and
$|\bar\lambda|=131100\leq md$.

These failures describe the scope of the comparison: the
example adds to the explicitly listed sufficient criteria.
It does not prove that every other conceivable method would
fail. The conclusion is the exact full determinant multiplicity
in (O-014.4), with at least four actual permanent copies.
The exact permanent multiplicity remains undetermined.
The example certifies strict refinement of the listed explicit
conditions; it is not an exhibited surviving $q>K_n$ candidate.
The later nested-support criteria F-016 and M-011 also handle
this example. O-015.c gives a separate four-copy example
beyond all the stabilization criteria, using M-011 directly.

\paragraph{Geometric dictionary.}
The bounded coefficient space is an affine linear subscheme
of the scheme of degree-$b$ forms, cut out by coefficient
coordinates outside the prescribed exponent box. Its coordinates
are selected by torus weights, using SGA~3, I, Proposition~4.7.3,
pp.~35--36. The two trees define a morphism from that affine
coefficient scheme to the scheme of linear matrices over
$\mathbb Z$, and the determinant identity proves factorization
through the padded-form morphism. F-015 then interprets
restriction as a torus-equivariant map of coordinate modules;
representations as comodules are SGA~3, I, Proposition~4.7.2,
pp.~34--35, and Schur functor naturality is Macdonald,
I Appendix~A, (5.3)--(5.6), pp.~155--156. The integral
factorization does not itself split irreducibles integrally:
the full-copy comparison is made after base change to
$\mathbb C$, with the explicit factorial invertibility in
O-012.b. Multiplication of the Hessian sections takes place
in the actual ambient coordinate algebra, and the positivity
claim comes from their nonzero restriction to the permanent
orbit rather than a formal identity in its $K_0$ class.
